\subsection{Angular integration of the vertex and self-energy terms}
\label{sec:comp-supp-vertex-angle}
The closed vertex matrices can now be interfered with the Born
amplitude and integrated over the elastic scattering angle.
Define $x=s-1$ and $y=1-u=x+t$, with $c$ the center-of-mass
photon scattering cosine. Since $t=-x^2(1-c)/(2s)$, the physical
limits and Jacobian are
\begin{equation}
 y_- =\frac{x}{1+x},\qquad y_+=x,\qquad
 \dd c=\frac{2s}{x^2}\dd y.
\end{equation}
Let $T_s$ and $T_u$ be the spin-averaged interferences from the two
vertex insertions and the self-energy insertion in each ordering,
with the common loop factor removed as above. We contract each
external photon pair with $-\eta_{\mu\nu}$. These diagram subsets
are not separately gauge invariant. Physical transverse polarization
sums can replace these contractions only after including the
spanning graphs, when the full Ward identities apply.

The traces use only the Clifford anticommutation relation and
on-shell momentum identities. They are recorded by
\texttt{compton\_virtual\_trace\_analytic.py}. Both vertex insertions
in a given ordering give the same interference. For the $s$ ordering,
\begin{equation}
 T_s(x,y)=\sum_f f(1+x)
       \left[\frac{a_f(x)}y+b_f(x)+c_f(x)y\right],
\end{equation}
where the symbols $f$ in the first column of
Table~\ref{tab:comp-vt-angle-coefficients} mean the moments and
self-energy coefficients already evaluated above. The entries
$\widehat c_Z$ and $1$ are constants, not functions of $1+x$.
\begin{table}[!htbp]
\centering
\begin{tabular}{c|ccc}
$f$ & $a_f$ & $b_f$ & $c_f$\\
\hline
$J_{00}$ & $16(x^2+2x+4)/x$ & $-8(x^3+2x^2+8x+8)/x^2$ & $0$\\
$J_{10}$ & $-8(x^2+2x+8)/x$ & $8(x^3+x^2+6x+8)/x^2$ & $16/x$\\
$J_{11}$ & $-8x$ & $8(x+2)$ & $-8/x$\\
$J_{20}$ & $-4(2x^2+7x+12)/x$ & $4(2x^2+13x+12)/x^2$ & $-4/x$\\
$J_{21}$ & $4(x-4)$ & $-4(2x^2+x-4)/x$ & $4(x+2)/x$\\
$\widehat a$ & $-4(x+2)^2/x^2$ & $2(3x^2+12x+8)/x^3$ & $2(x+1)/x^2$\\
$\widehat b$ & $-2(x^2+4x+8)/x^2$ & $16(x+1)/x^3$ & $2/x^2$\\
$\widehat c_Z$ & $-2(x+4)/x$ & $2(3x+4)/x^2$ & $2/x$\\
$1$ & $2(7x^2+28x+32)/x^2$ & $-2(x+4)(9x+8)/x^3$ & $-2(3x+4)/x^2$\\
\end{tabular}
\caption{Exact trace coefficients for the two $s$-channel vertex
insertions and the self-energy insertion together.}
\label{tab:comp-vt-angle-coefficients}
\end{table}
There is no $J_{22}$ contribution after the trace. The integral is
therefore elementary once the moments have been evaluated,
\begin{equation}
 I_s=\sum_f f(1+x)\left[a_f\ln(1+x)
       +b_f\frac{x^2}{1+x}
       +\frac{c_f x^2}{2}\left(1-\frac1{(1+x)^2}\right)\right].
 \label{eq:comp-vt-s-angle}
\end{equation}

For the crossed ordering the moment arguments depend on the
integration variable, $S=1-y$. Substituting the explicit moments
and self-energy coefficients reduces $T_u$ to rational functions
multiplying $1$, $\ln y$ and $\operatorname{Li}_2(1-y)$.
Partial fractions and integration by parts give the following
closed primitive. Define the real functions
\begin{align}
 H_{10}(y)&=\frac{\pi^2}6-\operatorname{Li}_2(1-y),\\
 H_{010}(y)&=2\operatorname{Li}_3(y+\ii0)
                        -\ln y\operatorname{Li}_2(y+\ii0).
\end{align}
Their derivatives are $H'_{10}=\ln y/(y-1)$ and
$H'_{010}=H_{10}/y$. For $y>1$, continuation from the upper
half-plane gives imaginary parts $\pi\ln y$ and
$\pi\ln^2y/2$ for the dilogarithm and trilogarithm respectively.
They cancel in $H_{010}$. Thus these expressions are real for all
$y>0$, including when the physical interval crosses $y=1$.

The primitive is
\begin{align}
 F_u(x,y)={}&-\frac{2(x^2+8)}x\ln^2y
       +\frac{8(x+1)}x H_{010}(y)
       +P(x,y)\ln y+Q(x,y)H_{10}(y)+R(x,y),\\
 P(x,y)={}&-\frac2{y-1}-\frac{4(x+8)}y
       +\frac{8x+12}{x}y
       +\frac{2\widehat c_Z(x^2-3x-4)+2x^2+10x}{x},\\
 Q(x,y)={}&-\frac{4x}{y^2}+\frac{16}y
                         +\frac{2x^2-18x+20}{x},\\
 R(x,y)={}&\frac{4(-2\widehat c_Z+x-8)}y
                 +\frac{2\widehat c_Z-12x-16}{x}y.
 \label{eq:comp-vt-u-primitive}
\end{align}
The apparent pole multiplying $\ln y$ at $y=1$ is removable,
since $\ln y/(y-1)\to1$. Direct differentiation recovers the
complete crossed-ordering interference, including its finite
on-shell subtraction. Consequently the angular integral of this
virtual subset is
\begin{equation}
 \boxed{\mathcal C_{\mathrm{vertex}+\Sigma}
 =\frac1{2x^2}\operatorname{Re}
       [I_s+F_u(x,x)-F_u(x,x/(1+x))].}
 \label{eq:comp-vt-integrated}
\end{equation}
This is an analytic contribution to the total cross section, with
no loop-parameter or angular integral left. It still contains
$\widehat c_Z=-4-4\ln\lambda$. Its regulator dependence is canceled
only in the complete sum with the spanning and real contributions.
